\section{Symmetric posterior experiments}
\label{sec:experiments}

The observations below are experiments about an input that is uniform on
$[4]$.  A tree combines such experiments in two ways: an observation passes
through one broadcast edge, and observations from different branches are
combined.  We record the exact posterior laws for these operations, keeping
separate the independent laws used in a calculation and the actual law of a
combined observation.  The normalized posterior update and its change of
observation law are standard in broadcasting recursions; see
\cite[Sections~2.2--2.4]{MezardMontanari2006}.

Write
\[
  \Deltafour=\left\{p\in[0,1]^4:\sum_{i=1}^4p_i=1\right\},
  \qquad \unif=(1/4,1/4,1/4,1/4),
\]
and let $\Sym$ be the group of permutations of $[4]$.
An experiment consists of a uniform input $I\in[4]$ and an observation $Y$.
Write $\kappa_i$ for the conditional law of $Y$ given $I=i$, and
$\kappa=\frac14\sum_i\kappa_i$ for its marginal law.  Since each
$\kappa_i$ is absolutely continuous with respect to $\kappa$, its posterior
vector $p(Y)\in\Deltafour$ satisfies
\begin{equation}
  p_i(y)=\frac14\frac{\mathrm d\kappa_i}{\mathrm d\kappa}(y),
  \qquad
  \frac{\mathrm d\kappa_i}{\mathrm d\kappa}(y)=4p_i(y).
  \label{eq:posterior-likelihood}
\end{equation}
These identities also define the posterior on a general observation space,
up to a $\kappa$-null set.  Denote by $\mu$ the law of $p(Y)$ under
$\kappa$.  In particular,
\[
  \int_{\Deltafour}p_i\,\mathrm d\mu(p)=\frac14
  \qquad(i\in[4]).
\]
Conversely, any probability law $\mu$ with this barycenter is realized by
taking the observation space to be $\Deltafour$ and its conditional laws
to be $4p_i\,\mathrm d\mu(p)$.
We call the experiment symmetric if $\mu$ is invariant under all
permutations of its four coordinates.  This is a condition on the law of
the posterior, not on an individual posterior vector.

The following coordinates remove the constant direction:
\begin{equation}
  L=
  \begin{pmatrix}
    1&1&-1&-1\\
    1&-1&1&-1\\
    1&-1&-1&1
  \end{pmatrix},
  \qquad x=Lp,
  \qquad A(p)=|Lp|^2.
  \label{eq:posterior-coordinates}
\end{equation}
Since $L^{\mathsf T}L=4\operatorname{Id}-\boldsymbol{1}
\boldsymbol{1}^{\mathsf T}$, where $\boldsymbol{1}=(1,1,1,1)^{\mathsf T}$,
\begin{equation}
  A(p)=4\sum_{i=1}^4(p_i-1/4)^2,
  \qquad 0\le A(p)\le3,
  \qquad
  \frac12\sum_{i=1}^4|p_i-1/4|\le\frac12\sqrt{A(p)}.
  \label{eq:quadratic-tv}
\end{equation}
Thus decay of the mean of $A$ will imply nonreconstruction.

\begin{lemma}[Bayes product]
\label{lem:bayes-product}
Let two experiments have the same uniform input and observations that are
independent conditional on that input.  Let $\mu$ and $\nu$ be their
individual posterior laws.  For $p,q\in\Deltafour$, define
\[
  U_i(p,q)=4p_iq_i,
  \qquad z(p,q)=\sum_{i=1}^4U_i(p,q),
  \qquad P_i(p,q)=\frac{U_i(p,q)}{z(p,q)}
  \quad\text{when }z(p,q)>0.
\]
Set $P(p,q)=\unif$ when $z(p,q)=0$.  The posterior law of the combined
experiment, denoted by $\mu\star\nu$, is characterized by
\begin{equation}
  \int_{\Deltafour} f(r)\,\mathrm d(\mu\star\nu)(r)
  =\int_{\Deltafour\times\Deltafour}
       z(p,q)f(P(p,q))\,\mathrm d\mu(p)\,\mathrm d\nu(q)
  \label{eq:product-law}
\end{equation}
for every bounded measurable $f$.  The measure on the right has total
mass one.  If both experiments are symmetric, their combined experiment
is symmetric.
\end{lemma}

\begin{proof}
Let $\kappa,\eta$ be the marginal observation laws, with posterior vectors
$p(y),q(w)$.  Conditional independence and
\eqref{eq:posterior-likelihood} show that the joint observation law has
density
\[
  \frac14\sum_{i=1}^4(4p_i(y))(4q_i(w))=z(p(y),q(w))
\]
with respect to $\kappa\otimes\eta$.  Bayes' formula then gives the
combined posterior $P$.  Pushing this identity forward to the two
posterior vectors proves \eqref{eq:product-law}.  Its normalization can
also be checked directly:
\[
  \int z(p,q)\,\mathrm d\mu(p)\,\mathrm d\nu(q)
  =4\sum_{i=1}^4
       \left(\int p_i\,\mathrm d\mu\right)
       \left(\int q_i\,\mathrm d\nu\right)=1.
\]
Finally, simultaneous coordinate permutation leaves $z$ unchanged and
permutes $P$ by the same permutation.  Invariance of $\mu$ and $\nu$
therefore implies invariance of $\mu\star\nu$.
\end{proof}

In \eqref{eq:product-law}, $p$ and $q$ are independent under
$\mu\otimes\nu$.  They generally are not independent under the actual
observation law $z\,\mathrm d\mu\,\mathrm d\nu$.  All factorizations of
expectations below use the former measure, before the factor $z$ is
applied.  The choice of $P$ at $z=0$ never affects a combined expectation.

\begin{lemma}[Passage through one edge]
\label{lem:channel}
Let $I$ be uniform on $[4]$, let $J$ be obtained from $I$ through the
channel $P_\lambda$, and let $Y$ be an experiment about $J$, independent
of $I$ conditional on $J$.  If $\mu$ is the posterior law for $J$ given
$Y$, then the posterior law for $I$ given $Y$ is the pushforward of
$\mu$ by
\begin{equation}
  C_\lambda(p)=\unif+\lambda(p-\unif),
  \qquad L C_\lambda(p)=\lambda Lp,
  \qquad A(C_\lambda(p))=\lambda^2A(p).
  \label{eq:channel-scaling}
\end{equation}
In particular, this operation preserves symmetry.
\end{lemma}

\begin{proof}
The channel is symmetric and doubly stochastic, so $J$ is uniform and
$\Pr(I=i\mid J=j)=P_\lambda(i\mid j)$.  If $p$ is the posterior of
$J$ given $Y$, the posterior of $I$ is consequently
\[
  \sum_{j=1}^4P_\lambda(i\mid j)p_j
  =\lambda p_i+\frac{1-\lambda}{4}.
\]
The marginal law of $Y$ is the same whether it is regarded as an
experiment about $I$ or about $J$.  This proves the pushforward statement.
The remaining identities follow from $L\unif=0$, and symmetry follows
because $C_\lambda$ commutes with coordinate permutations.
\end{proof}

We will also use observed mixtures of experiments.  Suppose an index $S$
is independent of the input and, conditional on $S=s$, the observation
has posterior law $\mu_s$.  If $S$ itself is observed, the resulting
posterior law is $\int\mu_s\,\mathrm d\Pr_S(s)$.  Hence symmetry and
any linear inequality on posterior laws are preserved under such
mixtures.  In the tree application, $S$ will be the offspring count.
The uninformative experiment has posterior law concentrated at $\unif$;
the experiment that reveals the input has posterior law assigning mass
$1/4$ to each of the four coordinate vectors.
